% Unofficial University of Cambridge Poster Template
% https://github.com/andiac/gemini-cam
% a fork of https://github.com/anishathalye/gemini
% also refer to https://github.com/k4rtik/uchicago-poster

\documentclass[final]{beamer}

% ====================
% Packages
% ====================

\usepackage[T1]{fontenc}
\usepackage{lmodern}
\usepackage[orientation=portrait,size=a1,scale=1.15]{beamerposter}
\usetheme{gemini}
\usecolortheme{nott}
\usepackage{graphicx}
\usepackage{booktabs}
\usepackage{tikz}
\usepackage{pgfplots}
\pgfplotsset{compat=1.14}
\usepackage{anyfontsize}

% ====================
% Lengths
% ====================

% If you have N columns, choose \sepwidth and \colwidth such that
% (N+1)*\sepwidth + N*\colwidth = \paperwidth
\newlength{\sepwidth}
\newlength{\colwidth}
\setlength{\sepwidth}{0.025\paperwidth}
\setlength{\colwidth}{0.45\paperwidth}

\newcommand{\separatorcolumn}{\begin{column}{\sepwidth}\end{column}}

% ====================
% Title
% ====================

\title{x86 Architecture}

\author{William Bland}

% ====================
% Footer
% ====================

% \footercontent{
%   % \href{https://github.com/TheBlandit/EMC-x86}{https://github.com/TheBlandit/EMC-x86} \hfill
%   % \href{https://exetermathematicsschool.ac.uk/}{Exeter Mathematics School}
% }

% ====================
% Logo
% ====================

\logoright{\includegraphics[height=3.5cm]{logos/ems_logo.png}\hspace{20ex}}
% \logoleft{\hspace{20ex}\includegraphics[height=3.0cm]{logos/x86.png}}

% ====================
% Body
% ====================

\begin{document}

% Refer to https://github.com/k4rtik/uchicago-poster
% logo: https://www.cam.ac.uk/brand-resources/about-the-logo/logo-downloads
% \addtobeamertemplate{headline}{}
% {
%     \begin{tikzpicture}[remember picture,overlay]
%       \node [anchor=north west, inner sep=3cm] at ([xshift=-2.5cm,yshift=1.75cm]current page.north west)
%       {\includegraphics[height=7cm]{logos/unott-logo.eps}}; 
%     \end{tikzpicture}
% }

\begin{frame}[t]
\begin{columns}[t]
\separatorcolumn

\begin{column}{\colwidth}
  \begin{block}{Overview of the x86 architecture}
    Every processor contains an architecture which describes how it executes machine code.
    A common architecture for desktop processors to use is x86, this means that the same programs can be run across different x86 processors without the need to be rewritten.

    Below is a brief history of x86 processors, summarising the major changes.
    \begin{itemize}
      \item The first x86 processor was the x8086 and was 16-bit, it used memory segmentation to allow it to separate the address space for processes, and to give it access to 1MiB of memory
      \item The i286 introduced memory protection which changed the segmentation model to prevent processes from accessing memory they did not own
      \item The i386 was the first 32-bit processor, introducing 32-bit registers and paging (which is another form of memory management in addition to segmentation)
      \item The i486 introduced an integrated x87 FPU (previously they were a separate processor)
      \item The intel Pentium 4 introduced 64-bit extensions
    \end{itemize}
  \end{block}

  \begin{block}{Memory Management}
    The kernel has to manage the usage of memory that processes use (give memory to programs that request it, and to prevent a program accessing the memory of another).

    In real mode, memory segmentation is used.
    When accessing memory, the value in the segment register is multiplied by 16 and added to the effective address, which produces the physical address of the memory.
    This allows a program to be split into logical sections (one for each segment register), allows access to 1MiB of memory (due to a 20-bit address bus), and the memory used by different processes are kept separate by keeping them in different segments.
    However it does not prevent a program of accessing memory it has not been allocated, and is difficult to manage memory spanning multiple segments.

    In protected mode, memory protection is enabled; this modifies the segmentation method so the segment registers are now segment selectors into a kernel-defined segment descriptor table.
    A segment descriptor contains the base address, limit (length of the segment), and flags (permissions) of the segment which provides greater flexibility of the use of memory segments.
    Unlike real mode, this prevents processes accessing memory they have not been allocated, however it still requires programs to manage their use of segments.

    Paging splits up the address space into 4KiB pages, and introduces the paging hierarchy which is a memory structure that the kernel manages and the hardware automatically traverses to convert a linear address into a physical address.
    This method provides greater control than segmentation by controlling the abilities of specific pages, and is completely transparent to the program (so is easier to use than segmentation).
    Furthermore paging allows consecutive linear addresses to use any physical address, it prevents memory segmentation and makes it easier for programs to allocate/free memory.
  \end{block}

  \begin{block}{Caching}
    Processors contain cache (which stores recently read/written data) to prevent having to access main memory (which is much slower than accessing cache).
    Cache is split up into cache lines, and x86 uses the MESI cache coherency protocol to keep the cache lines in different processors valid.
    Before writing to cache, the processor pushes writes to a store buffer, this means that a write may not be immediately globally visible, a memory fence can be used to flush the store buffer (which updates all the cache states); some instructions can be made atomic using the lock prefix, which locks the cache line for that instruction (this is used to implement mutexes).

    Furthermore there is another cache called the TLB (Translation Lookaside Buffer) which stores recent translations of the linear to physical address translation due to paging (which prevents unnecessary hardware assisted page walks).
    When updating a page table, the kernel must invalidate the TLB entrys which contain the invalid entrys. 
  \end{block}

  \begin{block}{Interrupts}
    Occasionally a hardware device will require immediate attention or an instruction could have raised an exception, this will require the current process to be interrupted and run code to handle the interrupt (which is called an Interrupt Service Routine (ISR)).
    Interrupt descriptors contain information on how the interrupt should be handled and are stored inside an interrupt descriptor table which its address is stored within the Interrupt Descriptor Table Register.
  \end{block}
\end{column}

\separatorcolumn

\begin{column}{\colwidth}
  \begin{block}{BIOS/UEFI}
    Vivamus congue volutpat elit non semper. Praesent molestie nec erat ac
    interdum. In quis suscipit erat. \textbf{Phasellus mauris felis, molestie
    ac pharetra quis}, tempus nec ante. Donec finibus ante vel purus mollis
    fermentum. 

    \begin{enumerate}
      \item \textbf{Morbi mauris purus}, egestas at vehicula et, convallis
        accumsan orci. Orci varius natoque penatibus et magnis dis parturient
        montes, nascetur ridiculus mus.
      \item \textbf{Cras vehicula blandit urna ut maximus}. Aliquam blandit nec
        massa ac sollicitudin. Curabitur cursus, metus nec imperdiet bibendum,
        velit lectus faucibus dolor, quis gravida metus mauris gravida turpis.
      \item \textbf{Vestibulum et massa diam}. Phasellus fermentum augue non
        nulla accumsan, non rhoncus lectus condimentum.
    \end{enumerate}
  \end{block}

  \begin{block}{Registers}
    The processor contains many registers which can be directly accessed using instructions.
    One type of register are general purpose registers which can be used by most instructions and so can be used for most operations, these registers are also split up into multiple registers, for example the 32-bit EAX register is located in the lower 32 bits of the 64-bit RAX register.
    Control registers control the operation of the processor, for example CR3 contains the base address of the paging hierarchy.
    Segment registers are used in segmentation and are always 16-bit.
    Model Specific Registers (MSRs) are only present on some processors and are accessed with specific instructions, one use of them is to enable long mode extensions.
  \end{block}

  \begin{block}{Operating modes}
    Etiam sit amet tempus lorem, aliquet condimentum velit. Donec et nibh
    consequat, sagittis ex eget, dictum orci. Etiam quis semper ante. Ut eu
    mauris purus. Proin nec consectetur ligula. Mauris pretium molestie
    ullamcorper. Integer nisi neque, aliquet et odio non, sagittis porta justo.

    \begin{itemize}
      \item \textbf{Sed consequat} id ante vel efficitur. Praesent congue massa
        sed est scelerisque, elementum mollis augue iaculis.
        \begin{itemize}
          \item In sed est finibus, vulputate
            nunc gravida, pulvinar lorem. In maximus nunc dolor, sed auctor eros
            porttitor quis.
          \item Fusce ornare dignissim nisi. Nam sit amet risus vel lacus
            tempor tincidunt eu a arcu.
          \item Donec rhoncus vestibulum erat, quis aliquam leo
            gravida egestas.
        \end{itemize}
      \item \textbf{Pellentesque facilisis dolor in leo} bibendum congue.
        Maecenas congue finibus justo, vitae eleifend urna facilisis at.
    \end{itemize}
  \end{block}

  \begin{block}{Syscalls (Linux)}
    A different kind of highlighted block.

    $$
    \int_{-\infty}^{\infty} e^{-x^2}\,dx = \sqrt{\pi}
    $$

    Interdum et malesuada fames $\{1, 4, 9, \ldots\}$ ac ante ipsum primis in
    faucibus. Cras eleifend dolor eu nulla suscipit suscipit. Sed lobortis non
    felis id vulputate.

    \heading{A heading inside a block}

    Praesent consectetur mi $x^2 + y^2$ metus, nec vestibulum justo viverra
    nec. Proin eget nulla pretium, egestas magna aliquam, mollis neque. Vivamus
    dictum $\mathbf{u}^\intercal\mathbf{v}$ sagittis odio, vel porta erat
    congue sed. Maecenas ut dolor quis arcu auctor porttitor.

    \heading{Another heading inside a block}

    Sed augue erat, scelerisque a purus ultricies, placerat porttitor neque.
    Donec $P(y \mid x)$ fermentum consectetur $\nabla_x P(y \mid x)$ sapien
    sagittis egestas. Duis eget leo euismod nunc viverra imperdiet nec id
    justo.
  \end{block}

  \begin{block}{Instructions}
  \end{block}
\end{column}
\separatorcolumn

\end{columns}
\end{frame}
\end{document}
